% ==========================================================================
% Paveikslas: LSTM atminties lastelės schema (siuolaikine, 3 vartu versija)
% Preambuleje reikia: \usepackage{tikz}  (ir \usepackage{float} del [H])
% Idejimas i teksta:  \input{paveikslas_lstm_lastele}
%
% PASTABA. Cia pavaizduota ta LSTM versija, kuri realizuota Keras/PyTorch:
% su uzmarsties vartais f_t (Gers, Schmidhuber, Cummins, 2000).
% Originalioje 1997 m. architekturoje uzmarsties vartu nera. Jei tekste
% aprasoma tik 1997 m. versija, is schemos reikia pasalinti:
%   1) eilute su  \node[vartai] (f) ...
%   2) eilute     \draw[valdymas] (f) -- (mf);
%   3) eilute su  \node[oper] (mf) ...
%   4) zyme $f_t$ ir uzrasa "uzmarsties vartai"
%   5) is \foreach saraso {2.0,3.9,5.3,7.6} isbraukti 2.0
% ==========================================================================

\begin{figure}[H]
    \centering
    \definecolor{ivestis}{RGB}{189,215,238}
    \definecolor{pasleptas}{RGB}{237,237,237}
    \definecolor{isvestis}{RGB}{198,224,180}
    \definecolor{vartaispalva}{RGB}{255,230,153}
    \begin{tikzpicture}[x=1cm, y=1cm,
        vartai/.style={circle, draw, fill=vartaispalva, minimum size=8mm,
                       inner sep=0pt, font=\scriptsize},
        kandidatas/.style={circle, draw, fill=ivestis, minimum size=8mm,
                       inner sep=0pt, font=\scriptsize},
        funkcija/.style={circle, draw, fill=white, minimum size=7mm,
                       inner sep=0pt, font=\scriptsize},
        oper/.style={circle, fill=white, minimum size=4.5mm,
                       inner sep=0pt, font=\large},
        srautas/.style={->, line width=0.5pt},
        valdymas/.style={->, gray!60, line width=0.4pt},
        zyme/.style={font=\scriptsize, inner sep=1.5pt},
        paaisk/.style={font=\tiny, text=gray!75, inner sep=1.5pt}]

        % lasteles kontūras
        \draw[rounded corners=3pt, draw=gray!70, fill=pasleptas!60]
            (0.6,-0.25) rectangle (9.0,4.9);

        % --- lasteles busenos linija (nuolatine klaidos karusele) ---
        \draw[srautas] (0.15,4.2) -- (9.55,4.2);
        \node[oper] (mf)   at (2.0,4.2) {$\odot$};
        \node[oper] (plus) at (4.6,4.2) {$\oplus$};
        \node[zyme, anchor=east] at (0.15,4.2) {$c_{t-1}$};
        \node[paaisk, anchor=east] at (0.15,4.55) {ankstesnė būsena};
        \node[zyme, anchor=west] at (9.55,4.2) {$c_{t}$};
        \node[zyme, anchor=south] at (7.8,4.45) {ląstelės būsena $c$};

        % --- vartai ir kandidatine busena ---
        \node[vartai]     (f) at (2.0,1.8) {$\sigma$};
        \node[vartai]     (i) at (3.9,1.8) {$\sigma$};
        \node[kandidatas] (g) at (5.3,1.8) {$\tanh$};
        \node[vartai]     (o) at (7.6,1.8) {$\sigma$};

        % kandidatines busenos ir ivesties vartu sandauga
        \node[oper] (mi) at (4.6,3.0) {$\odot$};
        \draw[valdymas] (i) -- (mi);
        \draw[valdymas] (g) -- (mi);
        \draw[srautas]  (mi) -- (plus);

        % uzmarsties vartai valdo lasteles busena
        \draw[valdymas] (f) -- (mf);

        % --- isvesties kelias ---
        \node[funkcija] (th) at (6.6,3.0) {$\tanh$};
        \draw[srautas] (6.6,4.2) -- (th);
        \draw[srautas] (th) -- (9.55,3.0);
        \node[oper] (mo) at (7.6,3.0) {$\odot$};
        \draw[valdymas] (o) -- (mo);
        \node[zyme, anchor=west] at (9.55,3.0) {$h_{t}$};

        % --- ivestis: h_{t-1} ir x_t ---
        \draw[line width=0.5pt] (0.15,0.95) -- (1.2,0.95);
        \draw[line width=0.5pt] (0.15,0.25) -- (1.2,0.25);
        \node[zyme, anchor=east] at (0.15,0.95) {$h_{t-1}$};
        \node[paaisk, anchor=east] at (0.15,1.3) {ankstesnė išvestis};
        \node[zyme, anchor=east] at (0.15,0.25) {$x_{t}$};
        \node[paaisk, anchor=east] at (0.15,-0.1) {sekos elementas};
        \draw[line width=0.5pt] (1.2,0.95) -- (1.2,0.25);
        \draw[line width=0.5pt] (1.2,0.6) -- (7.6,0.6);
        \fill (1.2,0.6) circle (1.2pt);
        \foreach \x in {2.0,3.9,5.3,7.6}
            \draw[srautas] (\x,0.6) -- (\x,1.4);

        % --- vartu zymejimai ---
        \node[zyme, anchor=west] at (2.15,2.9) {$f_{t}$};
        \node[zyme, anchor=east] at (3.75,2.7) {$i_{t}$};
        \node[zyme, anchor=west] at (5.45,2.7) {$\tilde{c}_{t}$};
        \node[zyme, anchor=west] at (7.75,2.5) {$o_{t}$};

        % --- uzrasai po paveikslu ---
        \node[zyme, anchor=north] at (2.0,-0.6) {užmaršties vartai};
        \node[zyme, anchor=north] at (4.6,-0.6) {įvesties vartai};
        \node[zyme, anchor=north] at (7.6,-0.6) {išvesties vartai};

    \end{tikzpicture}
    \caption{LSTM atminties ląstelė. Indeksas $t$ žymi elemento vietą sekoje. Viršutinė linija --
             ląstelės būsena, kuria informacija (ir atgalinio sklidimo klaidos
             signalas) keliauja per sekos žingsnius beveik nekeisdama mastelio:
             $c_t = f_t \odot c_{t-1} + i_t \odot \tilde{c}_t$. Vartai sprendžia,
             kiek sukauptos būsenos išsaugoti ($f_t$), kiek naujos informacijos į ją
             įrašyti ($i_t$) ir kiek jos atskleisti likusiam tinklui:
             $h_t = o_t \odot \tanh(c_t)$. Simbolis $\odot$ žymi daugybą, $\oplus$ -- sudėtį.}
    \label{fig:lstm-lastele}
\end{figure}
